\subsection{\problemblock{*Rail} Data Block}\label{sec:block-rail}

Rail launches and sled tests are special cases because rails constrain the vehicle to
move in only one direction. For these cases, a \problemblock{*rail} block is used
instead of \problemblock{*fly} blocks to determine body attitude and trajectory.
The rail or sled direction is taken from the trajectory initial conditions or from the
end of the preceding segment, and motion remains constrained in that direction for
the duration of the segment as described in \cref{sec:rail-launches}.

A rail launch is requested with

\begin{lstlisting}[style=taosmanual]
*rail launch cfstat=0.15 cfslid=0.01
\end{lstlisting}

where \statevar{cfstat} is the static coefficient of friction and
\statevar{cfslid} is the sliding, or kinematic, coefficient. Static friction is used
while velocity is less than 0.001~ft/s; above that speed the sliding coefficient is
used. Because of this switch, two segments are often used to simulate a launch: the
first covers motion from zero to 0.001~ft/s, and the second covers velocities above
0.001~ft/s. This represents the time required for thrust to overcome static friction
and begin moving the vehicle.

For a launch rail, TAOS assumes that the missile is attached so that it cannot move
backward. Negative acceleration in the rail direction is therefore ignored. Path length
\statevar{plength} can be used as the segment final condition to represent physical
rail length.

A sled test is requested with

\begin{lstlisting}[style=taosmanual]
*rail sled cfstat=0.17 cfslid=0.02
\end{lstlisting}

using the same friction definitions. Negative acceleration is allowed for a sled test
because the vehicle can accelerate to speed and then slow to a stop.
